% ===== Idioma y codificación =====
\usepackage[spanish, es-nodecimaldot]{babel}
\usepackage[utf8]{inputenc}
\usepackage[T1]{fontenc}
\usepackage{textcomp}

% ===== Tipografía moderna =====
\usepackage{helvet}
\renewcommand{\familydefault}{\sfdefault}
\usepackage{microtype}
\usepackage{xcolor}

% ===== Márgenes y formato =====
\usepackage[a4paper,top=2.8cm,bottom=2.8cm,left=3cm,right=3cm]{geometry}
\usepackage{titlesec}
\titleformat{\section}{\Large\bfseries}{\thesection.\ }{0.4em}{}
\titleformat{\subsection}{\large\bfseries}{\thesubsection.\ }{0.4em}{}
\titlespacing*{\section}{0pt}{1.6ex plus .5ex minus .2ex}{1.0ex}
\titlespacing*{\subsection}{0pt}{1.2ex plus .5ex minus .2ex}{0.7ex}

% ===== Colores =====
\definecolor{UMABlue}{RGB}{0,51,102}

% ===== Paquetes útiles =====
\usepackage{amsmath}
\usepackage{amssymb}
\usepackage{graphicx}
\usepackage{booktabs}
\usepackage{tabularx}
\usepackage{siunitx}
\usepackage{float}
\usepackage{enumitem}
\usepackage{listings}
\usepackage[hidelinks]{hyperref}
\usepackage{tikz}
\usetikzlibrary{arrows.meta, decorations.markings, patterns}
\sisetup{locale=DE}

% ===== Estilo para listados de código =====
\lstset{
  language=C++,
  basicstyle=\ttfamily\footnotesize,
  keywordstyle=\bfseries\color{UMABlue},
  commentstyle=\itshape\color{gray},
  stringstyle=\color{red!60!black},
  numbers=left,
  numberstyle=\tiny\color{gray},
  numbersep=6pt,
  frame=single,
  framerule=0.4pt,
  rulecolor=\color{gray!50},
  breaklines=true,
  tabsize=4,
  showstringspaces=false,
  captionpos=b,
  aboveskip=8pt,
  belowskip=6pt,
}
